\documentclass[12pt,a4paper]{article}
\usepackage{ugmath}
\usepackage{float}
\usepackage{placeins}
\usepackage [spanish]{babel}
\usepackage[labelfont=bf,labelsep=space]{caption}
\newcommand{\paren}[1]{\left( #1 \right)}
\newcommand{\alumno}{Ricardo León Martínez}
\newcommand{\materia}{Calculo Diferencial e Integral II}
\newcommand{\profesor}{Fernando Nuñez Medina}
\newcommand{\tarea}{Tarea 12}
\newcommand{\fecha}{16/5/2026}

\begin{document}
    \begin{center}
    {\large\textbf{UNIVERSIDAD DE GUANAJUATO}}\\[0.3cm]
    {\normalsize\textbf{DIVISIÓN DE CIENCIAS NATURALES Y EXACTAS}}\\
    {\normalsize\textbf{CAMPUS GUANAJUATO}}\\[1cm]

    {\Large\textbf{\tarea\ (\materia)}}\\[1cm]
    \end{center}

    \textbf{Nombre:} \alumno \hfill 
    \textbf{Fecha:} \fecha \hfill 
    \textbf{Calificación:} \rule{3cm}{0.4pt} \\[0.3cm]

    \begin{mdframed}[style=mdpinkbox,frametitle={Ejercicio 1}]
        Sean $(a_{n})$ y $(b_{n})$ sucesiones de números reales. ¿Se cumplirá que si $a_{n}\lt b_{n}$ para todo $n\gt N$,
        entonces $\limsup a_{n}\lt\limsup b_{n}$? Argumenta tu respuesta. Responde la misma
        pregunta para el limite inferior.
    \end{mdframed}

    \begin{proof}
        Las desigualdades estrictas no son ciertas en general. Consideremos el ejemplo
        \begin{equation*}
            a_{n}=0,\qquad b_{n}=\frac{1}{n}.
        \end{equation*}
        Entonces
        \begin{equation*}
            a_{n}\lt b_{n}\text{ para todo }n,
        \end{equation*}
        pero
        \begin{equation*}
            \limsup a_{n}=0,\qquad\limsup b_{n}=0.
        \end{equation*}
        Por tanto,
        \begin{equation*}
            \limsup a_{n}=\limsup b_{n},
        \end{equation*}
        y no se obtiene la desigualdad estricta. Además,
        \begin{equation*}
            \liminf a_{n}=0,\qquad \liminf b_{n}=0,
        \end{equation*}
        de modo que tampoco se cumple
        \begin{equation*}
            \liminf a_{n}\lt\liminf b_{n}.
        \end{equation*}
        En conclusión $a_{n}\lt b_{n}$ eventualmente implica únicamente
        \begin{equation*}
            \limsup a_{n}\leq\limsup b_{n},\qquad\liminf a_{n}\leq \liminf b_{n},
        \end{equation*}
        pero las desigualdades estrictas pueden fallar.
    \end{proof}

    \begin{mdframed}[style=mdpinkbox,frametitle={Ejercicio 2}]
        Determina la convergencia de las series siguientes:
        \begin{multicols}{2}
            \begin{enumerate}
                \item[(a)] $\sum_{n=1}^{\infty}\frac{\sqrt{n}}{n^{3}+9}$.
                \item[(b)] $\sum_{n=2}^{\infty}\frac{n}{\ln(n^{n})}$.
                \item[(c)] $\sum_{n=1}^{\infty}\frac{7}{n^{2}+3}$.
                \item[(d)] $\sum_{n=1}^{\infty}\frac{6}{1+4n}$.
            \end{enumerate}
        \end{multicols}
    \end{mdframed}

    \begin{solucion}
        \begin{enumerate}
            \item[(a)] $\sum_{n=1}^{\infty}\frac{\sqrt{n}}{n^{3}+9}$.\\
            Observemos que para $n$ grande, $n^{3}+9\sim n^{3}$. Por tanto,
            \begin{equation*}
                \frac{\sqrt{n}}{n^{3}+9}\sim\frac{\sqrt{n}}{n^{3}}=\frac{1}{n^{5/2}}.
            \end{equation*}
            Como
            \begin{equation*}
                \sum_{n=1}^{\infty}\frac{1}{n^{5/2}}
            \end{equation*}
            es una $p$-serie con
            \begin{equation*}
                p=\frac{5}{2}\gt1,
            \end{equation*}
            ésta converge. Aplicando el criterio de compración,
            \begin{equation*}
                \lim_{n\to\infty}\frac{\frac{\sqrt{n}}{n^{3}+9}}{\frac{1}{n^{5/2}}}=\lim_{n\to\infty}\frac{n^{3}}{n^{3}+9}=1.
            \end{equation*}
            Luego ambas series tiene el mismo carácter. Por tanto,
            \begin{equation*}
                \sum_{n=1}^{\infty}\frac{\sqrt{n}}{n^{3}+9}
            \end{equation*}
            converge.
            \item[(b)] $\sum_{n=2}^{\infty}\frac{n}{\ln(n^{n})}$.\\
            Usamos que
            \begin{equation*}
                \ln(n^{n})=n\ln(n).
            \end{equation*}
            Entonces
            \begin{equation*}
                \frac{n}{\ln(n^{n})}=\frac{n}{n\ln(n)}=\frac{1}{\ln(n)}.
            \end{equation*}
            La serie queda
            \begin{equation*}
                \sum_{n=2}^{\infty}\frac{1}{\ln(n)}.
            \end{equation*}
            Ahora,
            \begin{equation*}
                \lim_{n\to\infty}\frac{\frac{1}{\ln(n)}}{\frac{1}{n}}=\lim_{n\to\infty}\frac{n}{\ln(n)}=\infty,
            \end{equation*}
            de modo
            \begin{equation*}
                \frac{1}{\ln(n)}
            \end{equation*}
            es eventualmente mayor que una constante múltiplo de $1/n$. Como
            \begin{equation*}
                \sum\frac{1}{n}
            \end{equation*}
            diverge, por comparación se concluye que
            \begin{equation*}
                \sum_{n=2}^{\infty}\frac{1}{\ln(n)}
            \end{equation*}
            tambiém diverge. Por tanto,
            \begin{equation*}
                \sum_{n=2}^{\infty}\frac{n}{\ln(n^{n})}
            \end{equation*}
            también diverge.
            \item[(c)] $\sum_{n=1}^{\infty}\frac{7}{n^{2}+3}$.\\
            Para $n$ grande, $n^{2}+3\sim n^{2}$, y
            \begin{equation*}
                \frac{7}{n^{2}+3}\sim\frac{7}{n^{2}}.
            \end{equation*}
            Como
            \begin{equation*}
                \sum_{n=1}^{\infty}\frac{1}{n^{2}}
            \end{equation*}
            converge, usamos el criterio de compración,
            \begin{equation*}
                \lim_{n\to\infty}\frac{\frac{7}{n^{2}}}{\frac{7}{n^{2}}}=\lim_{n\to\infty}\frac{n^{2}}{n^{2}+3}=1.
            \end{equation*}
            Luego las dos tienen el mismo carácter. Por tanto
            \begin{equation*}
                \sum_{n=1}^{\infty}\frac{7}{n^{2}+3}
            \end{equation*}
            converge.
            \item[(d)] $\sum_{n=1}^{\infty}\frac{6}{1+4n}$.\\
            Veamos que para $n$ grande,
            \begin{equation*}
                \frac{6}{1+4n}\sim\frac{6}{4n}=\frac{3}{2n}.
            \end{equation*}
            Entonces comparamos con la serie armónica
            \begin{equation*}
                \sum\frac{1}{n}.
            \end{equation*}
            Usando el criterio de compración,
            \begin{equation*}
                \lim_{n\to\infty}\frac{\frac{6}{1+4n}}{\frac{1}{n}}=\lim_{n\to\infty}\frac{6n}{1+4n}=\frac{6}{4}=\frac{3}{2}.
            \end{equation*}
            El límite el positivo y finito, así que ambas series tienen el mismo carácter. Como
            \begin{equation*}
                \sum\frac{1}{n}
            \end{equation*}
            diverge, concluimos que
            \begin{equation*}
                \sum_{n=1}^{\infty}\frac{6}{1+4n}
            \end{equation*}
            diverge.
        \end{enumerate}
    \end{solucion}

    \begin{mdframed}[style=mdpinkbox,frametitle={Ejercicio 3}]
        Es claro que la expansión decimal de 0 es única y la de 1 no, pues $1=0.999\dots.$
        Sea $x\in(0,1)$. Muestra que la expansión decimal de $x$ es única ssi $x$ no tiene
        una expansión decimal finita y no tiene una expansión decimal de la forma $0.a_{1},\dots a_{n}\overline{9}$
        para algún numero natural $n$.
    \end{mdframed}

    \begin{proof}
        Probaremos la contraposición. Supongamos que $x$ tiene una expansión decimal finita.
        Entonces existe $n$ tal que
        \begin{equation*}
            x=0.a_{1}a_{2}\dots a_{n}.
        \end{equation*}
        Como
        \begin{equation*}
            0.999\dots=1,
        \end{equation*}
        se tiene
        \begin{equation*}
            0.a_{1}a_{2}\dots a_{n}=0.a_{1}a_{2}\dots(a_{n}-1)\overline{9},
        \end{equation*}
        si $a_{n}\neq0$. Por ejemplo,
        \begin{equation*}
            0.25000\dots=0.24999\dots.
        \end{equation*}
        Luego la expansión decimal no es única. Ahora supongamos que $x$ tiene una expansión
        de la forma
        \begin{equation*}
            x=0.a_{1},\dots a_{n}\overline{9}.
        \end{equation*}
        Entonces
        \begin{equation*}
            x=0.a_{1}a_{2}\dots(a_{n}+1),
        \end{equation*}
        si $a_{n}\neq9$. Por ejemplo,
        \begin{equation*}
            0.1399\dots=0.14.
        \end{equation*}
        Por tanto, nuevamente la expansión decimal no es única. Concluimos que si la expansión
        decimal de $x$ es única, entonces $x$ no puede tener expansión decirmal finita ni expansión
        terminada en nueves periodicos. Supongamos ahora que $x$ no tiene expansión decimal finita
        y tampoco tiene expansión decimal de la forma
        \begin{equation*}
            0.a_{1},\dots a_{n}\overline{9}.
        \end{equation*}
        Procederemos por contradicción. Supongamos que $x$ tiene dos expansiones decimales distintas
        \begin{equation*}
            x=0.a_{1}a_{2}a_{3}\dots
        \end{equation*}
        y 
        \begin{equation*}
            x=0.b_{1}b_{2}b_{3}\dots
        \end{equation*}
        Sea $m$ el menor número natural tal que $a_{m}\neq b_{m}$. Sin pérdida de generalidad,
        supongamos $a_{m}\lt b_{m}$. Entonces necesariamente
        \begin{equation*}
            a_{m}+1\leq b_{m}.
        \end{equation*}
        Usando la sugerencia, obtenemos la cadena de desigualdades
        \begin{gather*}
            x=0.a_{1}a_{2}a_{3}\dots a_{m-1}a_{m}a_{m+1}\dots\\
            \leq 0.a_{1}a_{2}a_{3}\dots a_{m-1}a_{m}\overline{9}\\
            =0.a_{1}a_{2}a_{3}\dots a_{m-1}(a_{m}+1)\\
            \leq0.a_{1}a_{2}a_{3}\dots a_{m-1}b_{m}\\
            =0.b_{1}b_{2}b_{3}\dots b_{m-1}b_{m}\\
            \leq0.b_{1}b_{2}b_{3}\dots b_{m-1}b_{m}b_{m+1}\dots=x.
        \end{gather*}
        Por hipótesis, $x$ no tiene expansión decimal terminada en nueves periodicos ni una expansión
        finita. Por ello, la primera y última desigualdad son estrictas. Así obtenemos
        \begin{equation*}
            x\lt x,
        \end{equation*}
        lo cual es un absurdo. La contradicción proviene de haber puesto que existían dos expansiones
        distintas. Así, la expansión decimal de $x$ es única.
    \end{proof}

    \begin{mdframed}[style=mdpinkbox,frametitle={Ejercicio 4}]
        Sea $x\in[0,1]$.
        \begin{enumerate}
            \item[(a)] Prueba que $x$ es racional ssi $x$ tiene una expansión decimal finita
            infinita pero periódica.
            \item[(b)] Prueba que $x$ es irracional ssi $x$ tiene una expansión decial infinita
            y no periodica.
        \end{enumerate}
    \end{mdframed}

    \begin{proof}
        \begin{enumerate}
            \item[(a)] Prueba que $x$ es racional ssi $x$ tiene una expansión decimal finita
            infinita pero periódica.\\
            Supongamos que $x$ es racional. Entonces existen enteros $p,q$ con
            \begin{equation*}
                q\neq0,\qquad x=\frac{p}{q}.
            \end{equation*}
            Al efectuar la división decimal de $p$ entre $q$, en cada paso el residuo posible
            pertenece al conjunto
            \begin{equation*}
                \{0,1,2,\dots,q-1\}.
            \end{equation*}
            Hay solamente un número finito de residuos posibles. Entonces  al realizar infinitamente
            muchos pasos ocurre una de las siguientes posibilidades: En algun momento aparece 0.
            En este caso la expansión decimal termina, es decir, es finita. Algun residuo se repite.
            Como el algoritmo de la división queda completamente determinado por el residuo, a partir
            de ese momento los dígitos comienzan a repetirse periódicamente. Por tanto, toda expansión
            decimal de un racional es finita o eventualmente periodica. Ahora supongamos que $x$ tiene
            expansión decimal finita o infinita periodica.\\
            Primero caso: expansión decimal finita\\
            Supongamos
            \begin{equation*}
                x=0.a_{1}a_{2}\dots a_{n}.
            \end{equation*}
            Entonces
            \begin{equation*}
                x=\frac{a_{1}a_{2}\dots a_{n}}{10^{n}},
            \end{equation*}
            que es racional.\\
            Segundo caso: expansión decimal periódica\\
            Supongamos
            \begin{equation*}
                x=0.a_{1}a_{2}\dots a_{k}\overline{b_{1}b_{2}\dots b_{m}}.
            \end{equation*}
            Escribamos $A+B$, donde
            \begin{equation*}
                A=0.a_{1}a_{2},\dots a_{k}\quad\text{ y }\quad B=0.00\dots0\overline{b_{1}b_{2}\dots b_{m}},
            \end{equation*}
            con $k$ ceros iniciales. Claramente $A$ es racional. Basta probar que $B$ es racional.
            Sea
            \begin{equation*}
                y=0.\overline{b_{1}b_{2}\dots b_{m}}
            \end{equation*}
            Entonces
            \begin{equation*}
                10^{m}y=b_{1}b_{2}\dots b_{m}\overline{b_{1}b_{2}\dots b_{m}}.
            \end{equation*}
            Restando,
            \begin{equation*}
                10^{m}y-y=b_{1}b_{2}\dots b_{m}.
            \end{equation*}
            Por tanto,
            \begin{equation*}
                (10^{m}-1)y=b_{1}b_{2}\dots b_{m}.
            \end{equation*}
            y así
            \begin{equation*}
                y=\frac{b_{1}b_{2}\dots b_{m}}{10^{m}-1}\in\Q.
            \end{equation*}
            Luego $B=\frac{y}{10^{k}}$ también es racional, y por lo tanto $x=A+B\in\Q$. Concluimos que
            $x$ es racional.
            \item[(b)] Prueba que $x$ es irracional ssi $x$ tiene una expansión decial infinita
            y no periodica.\\
            Supongamos que $x$ es irracional. Si su expansión decimal fuese finita o periódica, por
            el inciso (a) tendríamos que $x$ sería racional, contradicción. Por tanto, la expansión
            decimal de un irracional debe ser infinita no periódica. Supongamos ahora que $x$
            tiene expansión decimal infinita no periódica. Si $x$ fuese racional, entonces por el
            inciso (a) su expansión decimal tendría que ser finita o periódica, contradicción. Luego
            $x$ no es racional, es decir,
            \begin{equation*}
                x\in\R\setminus\Q.
            \end{equation*}
        \end{enumerate}
    \end{proof}

    \begin{mdframed}[style=mdpinkbox,frametitle={Ejercicio 5}]
        Prueba que si $f_{n}\rightrightarrows f$ en $D$ y $g_{n}\rightrightarrows f$ en $D$
        entonces $f_{n}+g_{n}\rightrightarrows f+g$ en $D$.
    \end{mdframed}

    \begin{proof}
        Por definición de convergencia uniforme, para todo $\e\gt0$ existe $N_{1}\in\N$ tal que
        \begin{equation*}
            n\geq N_{1}\implies \abs{f_{n}(x)-f(x)}\lt\frac{\e}{2}.
        \end{equation*}
        para toda $x\in D$. Asimismo, existe $N_{2}\in\N$ tal que
        \begin{equation*}
            n\geq N_{2}\implies\abs{g_{n}(x)-g(x)}\lt\frac{\e}{2}
        \end{equation*}
        para toda $x\in D$. Sea $N=\max\{N_{1},N_{2}\}$. Entonces, para todo $n\geq N$ y toda
        $x\in D$,
        \begin{equation*}
            \abs{f_{n}(x)-f(x)}\lt\frac{\e}{2},\qquad\abs{g_{n}(x)-g(x)}\lt\frac{\e}{2}.
        \end{equation*}
        Ahora,
        \begin{align*}
            \abs{(f_{n}+g_{n})(x)-(f+g)(x)}&=\abs{f_{n}(x)+g_{n}(x)-f(x)-g(x)}\\
            &=\abs{(f_{n}(x)-f(x))+(g_{n}(x)-g(x))}\\
            &\leq\abs{f_{n}(x)-f(x)}+\abs{g_{n}(x)-g(x)}\\
            &\lt\frac{\e}{2}+\frac{\e}{2}\\
            &=\e.
        \end{align*}
    \end{proof}

    \begin{mdframed}[style=mdpinkbox,frametitle={Ejercicio 6}]
        Prueba que si $f_{n}\rightrightarrows f$ en $D$, $g_{n}\rightrightarrows f$ en $D$
        y tanto $f_{n}$ como $g_{n}$ son acotadas en $D$ para todo $n\in\N$, entonces
        $f_{n}g_{n}\rightrightarrows fg$ en $D$.
    \end{mdframed}

    \begin{proof}
        Como $f_{n}\rightrightarrows f$, existe $N_{1}\in\N$ tal que
        \begin{equation*}
            n\geq N_{1}\implies \abs{f_{n}(x)-f(x)}\lt1
        \end{equation*}
        para toda $x\in D$. Entonces, para $n\geq N_{1}$,
        \begin{equation*}
            \abs{f(x)}\leq\abs{f(x)-f_{n}(x)}+\abs{f_{n}(x)}\lt1+\abs{f_{n}(x)}.
        \end{equation*}
        Como $f_{n}$ es acotada en $D$, existe $M_{n}\gt0$ tal que
        \begin{equation*}
            \abs{f_{n}(x)}\leq M_{n}\qquad\forall x\in D.
        \end{equation*}
        Por tanto,
        \begin{equation*}
            \abs{f(x)}\leq M_{n}+1,
        \end{equation*}
        de modo que $f$ también es acotada en $D$. Análogamente, $g$ es acotada. Tomemos
        constante $M,K\gt0$ tales que
        \begin{equation*}
            \abs{f(x)}\leq M,\qquad \abs{g_{n}(x)}\leq K
        \end{equation*}
        para toda $x\in D$ y todo $n$ suficientemente grande. Ahora escribimos
        \begin{equation*}
            f_{n}g_{n}-fg=f_{n}g_{n}-fg_{n}+fg_{n}-fg.
        \end{equation*}
        Entonces
        \begin{equation*}
            f_{n}g_{n}-fg=g_{n}(f_{n}-f)+f(g_{n}-g)-
        \end{equation*}
        Por la desigualdad del triangulo,
        \begin{equation*}
            \abs{f_{n}g_{n}-fg}\leq\abs{g_{n}}\abs{f_{n}-f}+\abs{f}\abs{g_{n}-g}
        \end{equation*}
        Por las acotaciones,
        \begin{equation*}
            \abs{f_{n}g_{n}-fg}\leq K\abs{f_{n}-f}+M\abs{g_{n}-g}.
        \end{equation*}
        Sea $\e\gt0$. Como $f_{n}\rightrightarrows f$, existe $N_{2}$ tal que
        \begin{equation*}
            n\geq N_{2}\implies\abs{f_{n}(x)-f(x)}\lt\frac{\e}{2K}
        \end{equation*}
        para toda $x\in D$. Asimismo, como $g_{n}\rightrightarrows g$, existe $N_{3}$
        tal que
        \begin{equation*}
            n\geq N_{3}\implies \abs{g_{n}(x)-g(x)}\lt\frac{\e}{2M}
        \end{equation*}
        para toda $x\in D$. Sea $N=\max{N_{1},N_{2},N_{3}}$. Entonces para toda $n\geq N$
        y para toda $x\in D$,
        \begin{align*}
            \abs{f_{n}(x)g_{n}(x)-f(x)g(x)}&\leq K\abs{f_{n}(x)-f(x)}+M\abs{g_{n}(x)-g(x)}\\
            &\lt K\frac{\e}{2K}+M\frac{\e}{2M}\\
            &=\e.
        \end{align*}
    \end{proof}

    \begin{mdframed}[style=mdpinkbox,frametitle={Ejercicio 7}]
        Calcula el radio de convergencia de las series de potencias siguientes:
        \begin{multicols}{2}
            \begin{enumerate}
                \item[(a)] $\sum_{n=0}^{\infty}x^{n}$.
                \item[(b)] $\sum_{n=0}^{\infty}nx^{n}$.
                \item[(c)] $\sum_{n=0}^{\infty}2^{n}x^{n}$.
                \item[(d)] $\sum_{n=0}^{\infty}n!x^{n}$.
            \end{enumerate}
        \end{multicols}
    \end{mdframed}

    \begin{solucion}
        \begin{enumerate}
            \item[(a)] $\sum_{n=0}^{\infty}x^{n}$\\
            Aquí $a_{n}=1$. Entonces por el teorema de Cauchy-Hadamard
            \begin{equation*}
                \sqrt[n]{\abs{a_{n}}}=\sqrt[n]{1}=1.
            \end{equation*}
            Por tanto, 
            \begin{equation*}
                \frac{1}{R}=1,
            \end{equation*}
            y así,
            \begin{equation*}
                R=1.
            \end{equation*}
            \item[(b)] $\sum_{n=0}^{\infty}nx^{n}$\\
            Aquí, $a_{n}=n$. Por el teorema de Cauchy-Hadamard
            \begin{equation*}
                \frac{1}{R}=\limsup_{n\to\infty}\sqrt[n]{n}.
            \end{equation*}
            Pero
            \begin{equation*}
                \lim_{n\to\infty}\sqrt[n]{n}=1.
            \end{equation*}
            Luego
            \begin{equation*}
                \frac{1}{R}=1, 
            \end{equation*}
            y por tanto $R=1$.
            \item[(c)] $\sum_{n=0}^{\infty}2^{n}x^{n}$\\
            Aquí, $a_{n}=2^{n}$. Entonces por el teorema de Cauchy-Hadamard
            \begin{equation*}
                \sqrt[n]{\abs{a_{n}}}=\sqrt[n]{2^{n}}=2.
            \end{equation*}
            Por tanto,
            \begin{equation*}
                \frac{1}{R}=2,
            \end{equation*}
            y así, $R=\frac{1}{2}$.
            \item[(d)] $\sum_{n=0}^{\infty}n!x^{n}$\\
            Aquí, $a_{n}=n!$. Por la proposición 102
            \begin{equation*}
                \abs{\frac{a_{n}}{a_{n+1}}}=\frac{n!}{(n+1)!}=\frac{1}{n+1}.
            \end{equation*}
            Tomando el limite
            \begin{equation*}
                \lim_{n\to\infty}\frac{1}{n+1}=0.
            \end{equation*}
            Por tanto, $R=0$. Así, la serie converge solamente en $x=0$.
        \end{enumerate}
    \end{solucion}
    
\end{document}